% Generated draft: requires mathematical and source review.
\documentclass[11pt]{article}
\usepackage[margin=25mm]{geometry}
\usepackage{fontspec}
\setmainfont{Latin Modern Roman}
\usepackage{amsmath,amssymb,mathtools,unicode-math}
\setmathfont{Latin Modern Math}
\usepackage{enumitem,needspace,xurl,hyperref}
\setlength{\parindent}{0pt}
\setlength{\parskip}{5pt}
\setlength{\emergencystretch}{3em}
\newcommand{\N}{\mathbb{N}}
\newcommand{\Z}{\mathbb{Z}}
\newcommand{\Q}{\mathbb{Q}}
\newcommand{\R}{\mathbb{R}}
\newcommand{\C}{\mathbb{C}}
\newcommand{\F}{\mathbb{F}}
\newcommand{\A}{\mathbb{A}}
\newcommand{\PP}{\mathbb{P}}
\newcommand{\E}{\mathbb{E}}
\newcommand{\eps}{\varepsilon}
\newcommand{\dd}{\,\mathrm{d}}
\begin{document}
\section*{UnsolvedMath ID: 11000031}
\textbf{Source problem code:} catalogue entry
\subsection{Definitions and mathematical statement}
Let $\operatorname{Mod}_g$ be the mapping class group in the source's genus-$g$ surface convention. For $H\leq\operatorname{Mod}_g$, let $\mathcal F(H)$ be its finite subgroups and identify $F,F'\in\mathcal F(H)$ when $F'=hFh^{-1}$ for some $h\in H$.

Problem 3.10 asks to find a finitely presented subgroup $H\leq\operatorname{Mod}_g$ with
\[
 \#\bigl(\mathcal F(H)/H\text{-conjugacy}\bigr)=\infty.
\]
Finite presentation means a presentation with finitely many generators and relators. Conjugacy is taken inside $H$, not in the ambient mapping class group; the problem is not asking for infinitely many isomorphism types of finite subgroups. The excerpt leaves the genus range and whether one seeks an example in some genus or every specified genus unstated.

The source motivates this with Bridson's example [Br] in $\operatorname{GL}_n(\Z)$. It suggests constructing an analogous subgroup in $\operatorname{Sp}_{2g}(\Z)$ and pulling it back under the homological action $\pi:\operatorname{Mod}_g\to\operatorname{Sp}_{2g}(\Z)$. Its phrase ``injective on torsion'' should be read with care: in the usual relevant range, torsion-freeness of the Torelli kernel implies injectivity on each finite subgroup, not necessarily injectivity on the entire set of torsion elements. The proposed strategy alone does not establish finite presentability of a pullback or the existence of suitable lifted finite subgroups.

The final extracted passage introduces a subsequent topic: the isomorphism problem for a collection $\mathcal S$ of finitely presented groups asks for an algorithm deciding, from two input presentations of members of $\mathcal S$, whether they define isomorphic groups. No additional isomorphism-problem demand is stated in this numbered item.
\subsection{Short English statement}
Find a finitely presented mapping-class subgroup with infinitely many conjugacy classes, within that subgroup, of finite subgroups. The suggested symplectic pullback strategy needs additional finiteness and lifting arguments.
\subsection{Sources}
\begin{itemize}
\item Local registry: \url{data/unsolvedmath/problems.json}, permanent UnsolvedMath ID $11000031$, exact source code \texttt{catalogue entry}.
\item Farb, \emph{Problems on Mapping Class Groups and Related Topics}, Problem 3.10, PDF page 34: \url{https://www.math.uchicago.edu/~farb/papers/mcgbook.pdf}. Supplied access date: 2026-07-29; extraction: pdf-font-structure; page-disambiguated numbering among 329 headings.
\item The source attributes the integral-linear example to Bridson [Br], with no full bibliography entry in this record. No precise modern citation supports the registry's report that no such mapping-class subgroup was verified through 2026.
\end{itemize}
\end{document}
